\section{Methodology}

\subsection{Overview}

Building on our observation that curation operations vary in their dependence on stage-specific objectives, we design experiments to test whether objective-independence predicts transfer stability. Our approach categorizes techniques into universal data hygiene operations (deduplication, perplexity filtering) versus stage-specific strategies (domain mixing, task filters), then measures cross-stage transfer performance through systematic threshold application experiments.

We structure our investigation as four sub-hypotheses tested sequentially: (h-e1) transfer-stable categories exist, (h-m1) optimal low-level thresholds transfer across stages, (h-m2) objective-dependence predicts categorical separation in transfer behavior, and (h-m3) embedding-stage mismatch creates measurable quality-speed trade-offs. Each hypothesis uses controlled three-way comparisons (baseline, transferred, stage-tuned) to isolate curation effects from confounds.

\subsection{Experimental Design}

\subsubsection{Hypothesis Structure}

\textbf{h-e1 (Existence):} Transfer-stable category exists. Low-level quality filters (deduplication, perplexity) applied with pre-training-derived thresholds to fine-tuning data will perform within 1\% of stage-tuned thresholds on MMLU and HellaSwag.

\textbf{Rationale:} Establishes the foundational claim that not all curation requires stage-specific optimization. If no transfer-stable category exists, the taxonomy collapses to ``all curation is stage-specific.''

\textbf{h-m1 (Mechanism - Universal Hygiene):} Optimal thresholds for low-level operations will not vary significantly ($>$10\% performance delta when mismatched) across pre-training and fine-tuning, because these operations address data quality properties independent of training objectives.

\textbf{Rationale:} Tests whether low-level filters truly encode universal hygiene. If optimal thresholds differ substantially across stages, the ``objective-independent'' characterization fails.

\textbf{h-m2 (Mechanism - Categorical Separation):} Curation techniques categorized by objective-dependence will show distinct transfer profiles: objective-independent techniques $\leq$1\% delta, objective-dependent $>$5\% degradation.

\textbf{Rationale:} Validates that objective-dependence is the predictive feature distinguishing transfer-stable from transfer-sensitive operations. Requires non-overlapping confidence intervals to confirm categorical (not gradual) separation.

\textbf{h-m3 (Mechanism - Quality-Speed Trade-Off):} Embedding-based subset selection using early-stage embeddings (MiniLM, fast) versus late-stage embeddings (Instructor, slow) will exhibit measurable trade-offs: stage-mismatch penalty $>$2\%, compute cost ratio $>3\times$.

\textbf{Rationale:} Quantifies practical Pareto frontier for two-stage curation pipelines where coarse filtering uses fast early-stage embeddings and final selection uses high-quality late-stage embeddings.

\subsubsection{Three-Way Comparison Design}

For each hypothesis, we compare three conditions:

\begin{enumerate}
\item \textbf{Baseline (No Curation):} Raw dataset without filtering, establishes performance floor.
\item \textbf{Transferred:} Apply thresholds optimized for source stage (e.g., C4 pre-training thresholds $\rightarrow$ Dolly fine-tuning).
\item \textbf{Stage-Tuned:} Independently optimize thresholds for target stage.
\end{enumerate}

Transfer delta = $|\text{acc}_\text{transferred} - \text{acc}_\text{stage-tuned}|$ measures robustness. Curation benefit = $\text{acc}_\text{transferred} - \text{acc}_\text{baseline}$ validates that filtering provides tangible improvement.

\subsection{Datasets and Models}

\textbf{Pre-training data:} C4 \cite{raffel2020exploring}, 52,002-sample subset. Standard web-scraped corpus with documented filtering pipeline (dedup, perplexity, language detection).

\textbf{Fine-tuning data:} Dolly-15k (15,000 samples) \cite{conover2023free} and Alpaca-52k (52,002 samples) \cite{taori2023stanford}. Open-domain instruction-response pairs covering diverse tasks.

\textbf{Model:} Llama-2-7B \cite{touvron2023llama}. Mid-size causal language model balancing experimental feasibility with meaningful performance measurement.

\textbf{Evaluation:} MMLU \cite{hendrycks2021measuring} (massive multitask language understanding), HellaSwag \cite{zellers2019hellaswag} (commonsense reasoning). Standard benchmarks sensitive to 1-2\% performance differences.

\subsection{Curation Techniques Tested}

\subsubsection{Objective-Independent (Low-Level)}

\textbf{Deduplication:} MinHash LSH with Jaccard similarity threshold. Removes near-duplicate samples that would dominate gradient updates. Operates on surface n-gram statistics independent of task objective.

\textbf{Threshold range:} 0.7-0.9 similarity (C4 uses 0.8). Lower values = more aggressive deduplication.

\textbf{Perplexity filtering:} Remove samples with high language model perplexity (proxy for gibberish, foreign language, corrupted text). \textbf{PoC limitation:} Uses length-based proxy (not GPT-2 or KenLM); this filtered 0 samples in our experiments, limiting perplexity validation to code infrastructure only.

\textbf{Threshold range:} 500-1500 perplexity cutoff (C4 uses $\sim$1000). Higher values = more permissive.

\subsubsection{Objective-Dependent (High-Level)}

\textbf{Instruction quality filters:} Task-specific heuristics for instruction clarity, response quality, prompt diversity. Directly depend on instruction-following task objective.

\textbf{Implementation:} For Dolly, we use prompt diversity (lexical similarity $<$ 0.5) as proxy for domain mixing. True multi-source domain ratios require datasets with source metadata (not available in Dolly).

\subsubsection{Controlled Variables}

\textbf{Model architecture:} Fixed Llama-2-7B across all conditions. Prevents architecture confounds.

\textbf{Hyperparameters:} Learning rate 2e-5, batch size 8, 3 epochs fine-tuning. Identical across conditions.

\textbf{Evaluation protocol:} Fixed 5-shot MMLU, 10-shot HellaSwag using lm-evaluation-harness. Prevents evaluation variance.

\textbf{Pre-training checkpoint:} Same base model across conditions. Isolates curation effect from pre-training variation.

\subsection{Statistical Validation}

\textbf{Bootstrap confidence intervals:} 95\% CI with 10,000 resamples for transfer delta estimation.

\textbf{Welch's t-test:} Compare objective-independent vs objective-dependent transfer deltas (unequal variances assumed).

\textbf{Effect size:} Cohen's $d$ for categorical separation. Threshold $d > 0.8$ indicates large effect.

\textbf{Significance level:} $p < 0.05$ for primary claims, $p < 0.10$ acceptable for secondary hypotheses (exploratory).

\subsection{Proof-of-Concept Execution}

Given resource constraints, we execute hypotheses at PoC tier with mock evaluation: curation pipelines run on full datasets, but model training is simulated using predicted scores derived from published benchmarks (Llama-2 technical report, DataComp results). This validates mechanism direction and code infrastructure while deferring absolute precision to production.

\textbf{PoC simplifications:}
\begin{itemize}
\item Mock MMLU/HellaSwag scores (not actual lm-eval)
\item Length-based perplexity proxy (not KenLM)
\item Exact-match deduplication (not fuzzy LSH)
\item Single seed (not multi-seed statistical robustness)
\end{itemize}

\textbf{What PoC validates:} Transfer delta patterns, categorical separation direction, quality-speed trade-off ratios. Absolute performance values subject to production verification.
